%%%%% Section 4. Geometrical Form: The Principal Connotative Space %%%%%
\section{Geometrical Form: The Principal Connotative Space}\label{sec:mathematics}

So far, denotation came first, and connotation lay beyond it. Now we reverse
the view. We begin with the nameless world of content data, and watch what a
name does there. One name is enough: a single facet. Its task here is to draw
out what description leaves behind. From it the space of the unnamed follows,
and we define that space in geometry, not in words.

%%% Subsection 4.1 %%%
\subsection{Content Space, Quale, and Qualia
    Spectrum}\label{subsec:content_space}

\paragraph{Ideal content space.}
As thermodynamics builds on the ideal gas, we build on the \textit{ideal
    content space} $\mathcal{X}$: an idealization of real content data. The ideal
gas averages how molecules move; $\mathcal{X}$ averages how people feel a
work. It is a continuous vector space, and every work that could be made
holds a unique position $x \in \mathcal{X}$.

A point of $\mathcal{X}$ carries the unique quality of one work, so we call a
point a \textit{quality}. Qualities also share aspects---jazz-ness, or a
quiet mood. We borrow a word for such a shared aspect: a \textit{quale}. A
quality is thus the whole of one work; a quale is one aspect running across
many works. Three axioms follow, and the third tells how a quale lies in
$\mathcal{X}$.

\begin{description}
    \item[\textbf{(Axiom 1) Sufficiency.}] The quality of a work carries what
          people feel in it---as far as the connotation we mean to reach.
    \item[\textbf{(Axiom 2) Similarity.}] $\mathcal{X}$ carries an inner
          product $\langle x, x' \rangle$, and closeness in it means likeness
          as people judge it, typically
          $\operatorname{Ave}\bigl[\operatorname{Sim}(w, w')\bigr] =
              \exp \langle x, x' \rangle$~\citep{Shepard19871317}.
    \item[\textbf{(Axiom 3) Manifold hypothesis.}] A quale is a probability
          density over $\mathcal{X}$, characterized by a smooth manifold.
\end{description}

\noindent
Axiom 3 turns sharing into a matter of degree. The density peaks on the manifold
and decays away from it: a work may carry a quale strongly, weakly, or hardly at
all. The manifold is the geometry of the quale, and we call it the \textit{quale
    manifold}---an idealized principal manifold, carrying the dominant structure of
the density and smooth enough to move along.

The three axioms together ask for one thing: a representation space faithful
to what people feel in a work. How that feeling is measured, and how such a
space is obtained, we leave open. We take it as given.

Current practice already reaches for such a space. The manifold hypothesis of
Axiom 3 has served machine learning for years~\citep{Bengio20131798}, and
modern embedding methods---Transformers and large language models---aim at
the very representation these axioms describe~\citep{Grand2022}.


\paragraph{Qualia spectrum.}
A quale can vary continuously---violet shades into blue, blue into green, and
so on to red---and its quale manifold moves with it. We call such a family of
manifolds a \textit{qualia spectrum}. Geometry already knows this structure as
a \textit{fiber bundle}, and we call it a \textit{spectrum fiber bundle}: the
base space $\mathcal{U}$ is the parameter---of any dimension---along which the
quale varies, and each fiber is a quale manifold.


\paragraph{Induced qualia spectrum.}
A spectrum fiber bundle induces a second spectrum. Fibers of one bundle share
a shape, so a single coordinate system $\mathcal{V}$ fits them all, and every
point of the bundle gains an address $(u, v)$.

Now hold $v$ fixed and let $u$ run, joining the point at $v$ on every fiber.
Those points form a manifold of their own, crossing the whole
bundle---geometry calls it a \textit{section}. Let $v$ vary, and these
sections sweep out a second spectrum: the \textit{induced qualia spectrum}.
The fibers are the warp, the sections the weft, and the qualities are the
cloth they weave.

How we cut $\mathcal{V}$ decides the cloth, and infinitely many cuts fit. To
choose one, we need an inductive bias. Ours has two parts.

\begin{description}
    \item[\textbf{(Bias 1) Independence.}] The density over the bundle
          factorizes: $p(u, v) = p(u)\,p(v)$.
    \item[\textbf{(Bias 2) Shortest weft.}] Among such coordinate systems,
          take the one whose sections run shortest in total.
\end{description}

\noindent
Bias 1 is what makes a quale common: $v$ says the same thing wherever it sits
on the base. Bias 2 makes it faithful: by Axiom 2, a short weft joins the
qualities that feel most alike, so each section is spun from like qualities
throughout. Together the two leave one bundle---disentangled by Bias 1, taut
by Bias 2. Shortest total length is the problem of optimal transport, so we
call it an \textit{optimal transport fiber bundle} (OT-FB), and what its
induced spectrum carries, the \textit{common qualia} of the base.

%%% Subsection 4.2 %%%
\subsection{Content Space Equipped with
    Denotation}\label{subsec:denotation}

\paragraph{Denotation and connotation.}
Section~\ref{subsec:naming} divided meaning by naming. That division now falls
on qualia. A quale that has a name is a \textit{denotation}; a quale that has
none is a \textit{connotation}. Their quale manifolds we call
\textit{denotative} and \textit{connotative}.

Naming leaves the geometry untouched. Jazz-ness and ``quiet, with feeling in
it'' are manifolds of one kind, and only one of them has a word. The
catalogue records the first; $\mathcal{X}$ holds both.

\paragraph{Facet.}
A qualia spectrum that has a name is a \textit{facet}. Genre is one: its
quale manifolds are jazz, classical, and the rest, and together they vary
with genre-ness. Facet analysis has long organized catalogues this
way~\citep{Cutter1876,Broughton200649}; here the facet becomes continuous,
and its parameter serves as the base coordinate $u$.

A named facet thus gives the base of a spectrum fiber bundle. What the
induced spectrum carries has no name---and that is where we go next.


%%% Subsection 4.3 %%%
\subsection{Facet-Induced Connotation Space}\label{subsec:principal_connotative}

\paragraph{Definition.}
A facet is the base of a spectrum fiber bundle, so the construction of
Section~\ref{subsec:content_space} applies unchanged. Its fibers are the
foci---jazz, classical, and the rest. Give them one coordinate system, take
the shortest weft, and the induced spectrum follows. Every work then sits at
an address $(u, v)$: $u$ runs across the foci, $v$ picks a section.

What does $v$ carry? A quale common to every focus, and one the facet cannot
name: genre names jazz and classical, nothing more. A quale without a name is
a connotation (Section~\ref{subsec:denotation}), so we call $\mathcal{V}$ the
\textit{facet-induced connotation space}.

\paragraph{Facet-induced OT-FB.}
The bundle so built we call the \textit{facet-induced OT-FB}. It is the
multi-floor library with rabbit holes itself, now in geometry
(Figure~\ref{fig:two_views}): the foci are its floors, the sections its
rabbit holes. Section~\ref{sec:framework} postulated such a library; here it
exists.

And it is unique. Give a collection and one facet, and the two biases leave
one bundle, essentially unique under mild regularity. Nothing more is asked
of the analyst. Appendix~\ref{sec:computation} gives the proof and the
algorithm.